\documentclass[12pt,a4paper]{article}
\usepackage[T1]{fontenc}\usepackage[utf8]{inputenc}\usepackage{times}
\usepackage[a4paper,top=2.5cm,bottom=2.5cm,left=2.5cm,right=2.5cm]{geometry}
\usepackage{graphicx}
\graphicspath{{figures/}}
\usepackage{booktabs}\usepackage{array}\usepackage{enumitem}\usepackage{tabularx}
\usepackage{caption}\usepackage{fancyhdr}\usepackage{parskip}\usepackage{float}
\usepackage{setspace}\onehalfspacing
\usepackage{amsmath}
\usepackage{xurl}\usepackage[hidelinks]{hyperref}
\captionsetup{font=small,labelfont=bf,labelsep=colon,
              justification=justified,singlelinecheck=false}
\setlist[itemize]{leftmargin=1.3em,itemsep=0pt,parsep=0pt,topsep=2pt}
\setlist[enumerate]{leftmargin=1.6em,itemsep=0pt,parsep=0pt,topsep=2pt}
\setlength{\parskip}{4pt plus 1pt}
\renewcommand{\topfraction}{0.85}\renewcommand{\bottomfraction}{0.7}\renewcommand{\textfraction}{0.1}\renewcommand{\floatpagefraction}{0.75}
\pagestyle{fancy}\fancyhf{}
\fancyhead[L]{\footnotesize KOUAME Koffi Fid\`ele}
\fancyhead[C]{\footnotesize VIDEO GAME MARKET ANALYTICS}
\fancyhead[R]{\footnotesize Data Analysis Internship}
\fancyfoot[C]{\thepage}\renewcommand{\headrulewidth}{0.4pt}
\setlength{\headheight}{15pt}
\newcommand{\fig}[3][0.92]{\begin{figure}[H]\centering\includegraphics[width=#1\linewidth]{#2}\caption{#3}\end{figure}}
\newcommand{\capture}[2]{\IfFileExists{figures/#1}{\begin{figure}[H]\centering\setlength{\fboxsep}{0pt}\setlength{\fboxrule}{0.4pt}%
\fbox{\includegraphics[width=\dimexpr\linewidth-0.8pt\relax]{#1}}\caption{#2}\end{figure}}{}}

\begin{document}

\begin{titlepage}\centering
\singlespacing
\vspace*{3cm}
{\large\scshape Data Analysis Internship}\\[0.4cm]
{\normalsize Task 11}\\[2.0cm]
\rule{\linewidth}{0.5pt}\\[0.8cm]
{\LARGE\bfseries Video Game Industry Analysis}\\[0.6cm]
{\large What Sells in the Video Game Market?\\ Genres, Platforms, Publishers, Regions and Ratings\\ of 16{,}715 Releases (1980--2016)}\\[0.8cm]
\rule{\linewidth}{0.5pt}\\[1.2cm]
{\normalsize Data Quality $\bullet$ Cleaning $\bullet$ Exploratory Analysis $\bullet$ Statistical Tests
$\bullet$ Power BI Dashboard $\bullet$ Recommendations}\\[3.0cm]
{\large\itshape Prepared by}\\[0.5cm]
{\Large\bfseries KOUAME Koffi Fid\`ele}\\[0.6cm]
{\large 5 October 2026}
\vfill\end{titlepage}

\singlespacing
\tableofcontents\newpage
\onehalfspacing

% =====================================================================
\section{Executive summary}

A company of the gaming industry wants to understand the video game market. This report
analyses a dataset of game sales to answer seven questions: which genres are the most
common, which platforms and publishers sell the most, which games are the best sellers,
how sales have changed over time, which regions buy, and whether ratings go with sales.
The notebook \texttt{Video-Game-Analytics.ipynb} reproduces every number and figure; the
cleaned dataset is in \texttt{data/processed/}, and the Power BI dashboard in
\texttt{powerbi/}.

\textbf{Data.} The main file lists \textbf{16{,}719 releases} (one game on one platform) on
31 platforms from 1980 to 2016, with sales in four regions, critic and user scores and the
ESRB age rating (VGChartz sales, Metacritic scores, extraction of 22 December 2016). Two
additional files give a later snapshot of the PS4 and Xbox One catalogues. Sales are
\textbf{millions of units}, not dollars. After cleaning, \textbf{16{,}715 releases of
11{,}563 games} remain, with 99.97\% of the units.

\textbf{Key figures.} \textbf{8{,}918 million units} sold. \textbf{Action} is both the most
common and the best-selling genre, the \textbf{PlayStation~2} the best-selling platform
and \textbf{Nintendo} the leading publisher. The typical release sells 170{,}000 copies;
only 12.4\% pass one million.

\textbf{Main findings.}
\begin{itemize}
\item \textbf{A hit-driven market}: 6.6\% of the releases make half of all units.
\item \textbf{Frequency is not performance}: Action is the most common genre, but
Platform and Shooter games sell the most per release; Adventure is crowded with small sellers.
\item \textbf{Nintendo} sells 20\% of all units with 2.5 million per release; the ten
largest publishers hold 70\% of the market.
\item \textbf{Retail sales peaked in 2008} (672 million units) and then fell by 60\%, as
digital and mobile games, which are not tracked, grew.
\item \textbf{North America buys half of the units}; Europe has risen to the first rank
for 2016 releases; Japan is a distinct market that prefers role-playing games.
\item \textbf{Critic scores go with sales} (Spearman $\rho = 0.39$), \textbf{user scores
much less} ($\rho = 0.15$); the median release scored 90+ sells 13 times more than one
scored below 50.
\item \textbf{Games sell for years}: the same PS4 and Xbox One games sold 29\% more after
December 2016.
\end{itemize}

\textbf{Dashboard.} The findings are delivered in an eight-page interactive Power BI dashboard
(Section~\ref{sec:dash}, Figure~\ref{fig:p1}): an overview with the KPI cards required by the brief,
one page per group of questions, an update on the PS4 and Xbox One generation, the insights and
an auditable data-quality page.

% =====================================================================
\section{Data understanding}

\subsection{The three files}

\begin{table}[H]\centering\small
\caption{Supplied files.}
\begin{tabularx}{\linewidth}{@{}lXr@{}}
\toprule
File & Content & Rows\\\midrule
\texttt{Video\_Games\_Sales\_as\_at\_22\_Dec\_2016.csv} & main file: one row per release; name, platform, year, genre, publisher, sales in NA, EU, JP, other regions and global, critic score and count, user score and count, developer, ESRB rating & 16{,}719\\
\texttt{PS4\_GamesSales.csv} & later VGChartz snapshot of PS4 games: game, year, genre, publisher, sales by region & 1{,}034\\
\texttt{XboxOne\_GameSales.csv} & same for Xbox One (plus a rank column) & 613\\
\bottomrule
\end{tabularx}
\end{table}

Two notions must be kept apart: a \textbf{release} is a game on one platform (one row,
16{,}715 after cleaning) and a \textbf{game} is a title, all platforms together (11{,}563).
\emph{Grand Theft Auto V}, for example, is one game with five releases.

\subsection{Columns, meaning and types}

\begin{table}[H]\centering\small
\caption{Columns of the main file, with the type read by pandas and the missing values.}
\begin{tabularx}{\linewidth}{@{}lXllr@{}}
\toprule
Column & Meaning & Read as & Correct type & Missing\\\midrule
Name & title of the game & text & text & 2\\
Platform & console or system code (PS2, X360, Wii, PC...) & text & text & 0\\
Year\_of\_Release & year of release on that platform & decimal & integer & 269\\
Genre & 12 genres & text & text & 2\\
Publisher & company that published the release & text & text & 54\\
NA/EU/JP/Other\_Sales & units sold by region (millions) & decimal & decimal & 0\\
Global\_Sales & units sold worldwide (millions) & decimal & decimal & 0\\
Critic\_Score & Metacritic press score (0--100) & decimal & decimal & 8{,}582\\
Critic\_Count & number of critics & decimal & integer & 8{,}582\\
User\_Score & Metacritic player score (0--10) & decimal & decimal & 9{,}129\\
User\_Count & number of players & decimal & integer & 9{,}129\\
Developer & studio & text & text & 6{,}623\\
Rating & ESRB age rating & text & text & 6{,}769\\
\bottomrule
\end{tabularx}
\end{table}

Sales are extremely skewed: median 0.17, mean 0.53 and maximum 82.5 million units
(\emph{Wii Sports}). Critic and user scores exist for less than half of the releases,
almost none before 1996: this is a coverage limit of Metacritic, not an error to correct.

% =====================================================================
\section{Data quality and cleaning}

\subsection{Issues found}

\begin{table}[H]\centering\small
\caption{Data quality audit and treatment.}
\begin{tabularx}{\linewidth}{@{}Xrl@{}}
\toprule
Issue & Rows & Treatment\\\midrule
No name and no genre (Genesis, 1993) & 2 & removed (2.42 M units)\\
Split record of one release: \emph{Madden NFL 13} (PS3), \emph{Sonic the Hedgehog} (PS3) & 2 & merged (sales added)\\
Two different games with one name: \emph{Need for Speed: Most Wanted} 2005 and 2012 & 2 pairs & kept, 2012 releases renamed\\
Year after the extraction date (2017, 2020) & 4 & year set to missing\\
Missing year & 269 & kept, \emph{Unknown}, out of time series\\
Missing publisher & 54 & \emph{Unknown}\\
Legacy ESRB code K-A (= E before 1998), pending code RP & 3 + 3 & K-A $\to$ E; RP and blank $\to$ \emph{Not rated}\\
Years and counts stored as decimals & 3 columns & integers\\
Exact duplicates; negative or zero sales & 0 & none found\\
Console files: Mac Roman encoding, carriage-return line breaks, non-breaking spaces, 512 zero-sales rows (announced games) & 1{,}647 & decoded, trimmed, zero-sales rows removed\\
\bottomrule
\end{tabularx}
\end{table}

The duplicate check shows why a mechanical rule is not enough: three name-platform pairs
appear twice, but two of them are fragments of one release (same critic score, tiny or
partial sales, missing year) and must be merged, while the third is a 2012 reboot of a
2005 game and must be kept. The 2020 release of \emph{Imagine: Makeup Artist} on DS, a
console discontinued in 2013, is a typing error: its year is set to missing rather than
guessed. The console files, read with default settings, appear empty, because they use
the old Mac line breaks; they are decoded with the Mac Roman encoding (byte 142 =
\emph{é}).

\subsection{Cleaned dataset}

Every step is recorded in a cleaning log (\texttt{data/processed/cleaning\_log.csv}),
also shown in the dashboard: \textbf{16{,}719 rows $\to$ 16{,}715 releases}, units
8{,}920.30 $\to$ \textbf{8{,}917.88 million (99.97\%)}. The cleaned dataset
\texttt{video\_games\_clean.csv} has 25 columns: the original ones, plus platform name,
manufacturer (Sony, Nintendo, Microsoft, Sega, PC, Other) and type (home console,
handheld, PC), decade, user score $\times$10 (same scale as the critic score), critic
score band and sales tier. Three companion tables feed the dashboard: the regional sales
in long format (66{,}860 rows), the cleaned console files (1{,}135 releases with sales)
and the cleaning log.

% =====================================================================
\section{Key performance indicators}

\begin{table}[H]\centering\small
\caption{KPIs of the dashboard (no filter).}
\begin{tabularx}{\linewidth}{@{}lrX@{}}
\toprule
KPI & Value & Comment\\\midrule
Total games & 11{,}563 & distinct titles\\
Total releases & 16{,}715 & game $\times$ platform\\
Total sales & 8{,}917.9 M units & 1980--2016\\
Top genre & Action & also the most common (3{,}370 releases)\\
Top platform & PlayStation 2 & 1{,}256 M units\\
Top publisher & Nintendo & 1{,}789 M units\\
Average / median sales per release & 0.53 / 0.17 M & skewed by blockbusters\\
Releases above 1 M units & 12.4\% & \\
\bottomrule
\end{tabularx}
\end{table}

% =====================================================================
\section{Exploratory analysis}

\subsection{Which genres are the most common? (Q1)}
\fig{fig01_genres}{Releases per genre (left) and units per release (right; grey: mean, blue: median).}
\textbf{Action is the most common genre} (3{,}370 releases, 20\%), ahead of Sports (14\%)
and Misc (10\%). Frequency and performance differ: \textbf{Platform} games are only the 8th
most common genre but sell the most per release (mean 0.93 M, median 0.27 M), followed by
\textbf{Shooter} (0.80 M). \textbf{Adventure} is the 6th most common genre but the weakest
seller (median 0.05 M); Strategy and Puzzle are similar. In every genre the median is
three to four times below the mean: each genre lives on a few hits.

\subsection{Which platforms have the highest sales? (Q2)}
\fig{fig02_platforms}{Global sales of the 15 best-selling platforms, coloured by manufacturer.}
The \textbf{PlayStation 2} leads (1{,}256 M units, 14.1\%), ahead of the Xbox 360 (972 M),
PS3 (939 M), Wii (908 M) and DS (807 M), all consoles of the 2000s, the peak of retail
sales. By manufacturer, \textbf{Sony (40\%) and Nintendo (39\%)} share four fifths of the
market, Microsoft 16\%. The PS4, launched in 2013, already ranks 8th: it leads the most
recent period.

\subsection{Which publishers have the highest total sales? (Q3)}
\fig{fig03_publishers}{Global sales of the 12 largest publishers, with their number of releases and units per release.}
\textbf{Nintendo} leads with 1{,}789 M units (20\%), followed by Electronic Arts (1{,}117 M),
Activision (731 M), Sony Computer Entertainment (606 M) and Ubisoft (472 M). Two models
appear: Nintendo publishes few releases (706) that sell 2.53 M each, five times the
average, on its own consoles; EA publishes almost twice as many (1{,}355) at 0.82 M each,
with annual franchises on every platform. The market is concentrated: the 10 largest of
581 publishers hold \textbf{70\%} of the units.

\subsection{What are the best-selling games? (Q4)}
\fig{fig04_best_sellers}{Ten best-selling releases (left) and ten best-selling games, all platforms together (right).}
\textbf{Wii Sports} (82.5 M) is by far the best-selling release, and the ten best-selling
releases are all Nintendo games on Nintendo hardware, several of them bundled with the
console. Summing the platforms of each title changes the ranking: \textbf{Grand Theft Auto~V}
becomes the second best-selling game (56.6 M on 5 platforms), and two \emph{Call of Duty}
episodes (about 31 M each) enter the top~10. For a third-party publisher, the route to the
top is the multi-platform blockbuster.

\subsection{How have game sales changed over the years? (Q5)}
\fig{fig05_sales_over_years}{Units sold (top) and releases (bottom) per year of release.}
The market went through three phases: a small market until 1995 (under 100 M units a year);
fast growth from 1996 to the \textbf{peak of 2008} (672 M units, 1{,}427 releases); then a
decline of tracked retail sales, to 268 M units in 2015 ($-60\%$). Releases and units fell
at the same pace, so each release sells about as much as before. The decline concerns
\emph{retail units}: digital downloads, mobile and free-to-play games are not tracked, and
2016 is incomplete.
\fig[0.92]{fig06_manufacturers_over_years}{Units sold per year by platform manufacturer.}
Sales follow \textbf{console generations} of six to eight years: Sony leads with the
PlayStation and PS2, Nintendo peaks with the Wii and DS (2006--2010), Microsoft grows
with the Xbox 360; Sega leaves hardware after 2001.

\subsection{Which regions generate the highest sales? (Q6)}
\fig{fig07_regions}{Share of units by region: all years (left) and per year (right).}
\textbf{North America generates half of all units} (49.4\%), ahead of Europe (27.2\%),
Japan (14.6\%) and the rest of the world (8.9\%). Japan weighed 27--29\% in the 1980s and
1990s but 10--15\% since 2005; Europe rose from 8\% (1980s) to 33\% (2010s) and is the
first region for 2016 releases (39\% vs 35\%).
\fig[0.9]{fig08_genre_by_region}{Share of each region's units by genre (each column adds up to 100\%); the red frames mark the two cells where Japan differs most.}
\textbf{Japan is a different market}: role-playing games make 27.4\% of its units (7.5\%
in North America), shooters only 3.0\% (13.5\%). North America, Europe and the rest of the
world buy the same mix of Action, Sports and Shooter games.

\subsection{Is there a relationship between ratings and sales? (Q7)}
The analysis uses the 8{,}135 releases with a critic score (48.7\% of the releases but
62.9\% of the units: reviewed games are the more visible ones).
\fig{fig09_critic_score_sales}{Critic score and units sold: every release (left, log scale) and median per score band (right).}
\textbf{Better-reviewed games sell more, and the effect accelerates at the top}: the
median release sells 0.12 M units below a score of 50, 0.28 M at 70--79, 0.59 M at 80--89
and \textbf{1.54 M at 90 or more}, thirteen times more than the lowest band. The cloud is
wide, however: the score shifts the odds, it does not guarantee a hit.
\fig{fig10_scores_esrb}{Rank correlations with sales (left) and median units by ESRB rating (right).}
\textbf{The press matters more than the players}: $\rho = 0.39$ for the critic score,
$0.15$ for the user score, although the two scores agree ($\rho = 0.54$). Critic reviews
appear at launch, when most copies are sold; user scores accumulate later. Mature (M)
games have the highest median (0.31 M), unrated releases the lowest (0.11 M).

\subsection{A hit-driven market}
\fig[0.86]{fig11_concentration}{Cumulative share of units, releases sorted from best to worst seller.}
\textbf{6.6\% of the releases make half of all units} and one quarter makes 80\%, while
35\% of the releases sell fewer than 100{,}000 copies.

% =====================================================================
\section{Statistical tests}

Sales are very skewed, so the tests are rank-based. With thousands of releases even weak
effects are significant; the size of the effect is what matters.

\begin{table}[H]\centering\small
\caption{Spearman correlations with global sales (95\% bootstrap intervals) and Kruskal-Wallis tests.}
\begin{tabularx}{\linewidth}{@{}Xrrl@{}}
\toprule
Relationship & $n$ & Effect & Reading\\\midrule
Critic score -- sales & 8{,}135 & $\rho = 0.39$ [0.37, 0.41] & moderate, strongest\\
Critic score -- sales, controlling for the number of critics & 8{,}135 & partial $\rho = 0.26$ & survives the visibility effect\\
Critic score -- sales within each genre & 12 genres & $\rho$ 0.13 to 0.50 & positive in all genres\\
Number of critics -- sales & 8{,}135 & $\rho = 0.41$ & big games get more reviews\\
User score -- sales & 7{,}588 & $\rho = 0.15$ [0.13, 0.17] & weak\\
Year of release -- sales & 16{,}443 & $\rho = -0.16$ & recent releases sell slightly less each\\
Sales by critic band (Kruskal-Wallis) & 8{,}135 & $\varepsilon^2 = 0.16$ & large\\
Sales by genre / ESRB / manufacturer & & $\varepsilon^2$ 0.06 / 0.04 / 0.03 & small\\
\bottomrule
\end{tabularx}
\end{table}

The critic score is the best single predictor of sales in this file, and the link holds
once the number of critics (a proxy of visibility) is removed. These are associations,
not causes: a large budget can produce both a better game and a larger marketing campaign.

% =====================================================================
\section{After 2016: the PS4 and Xbox One generation}

\fig{fig12_current_generation}{Region shares (left) and sales of the same 620 PS4 and Xbox One games in December 2016 and in the later snapshot (right).}
The console files are not merged with the main file (other snapshot, other genre names);
they answer two questions. \textbf{In the current generation, Japan almost disappears}
(4\% of PS4 and Xbox One units) and \textbf{Europe (39\%) catches up with North America
(43\%)}; the leaders are Western blockbusters (GTA~V, Call of Duty, Red Dead Redemption~2,
FIFA). \textbf{Games keep selling}: the 620 releases found in both files sold 453 M units
by December 2016 and 585 M in the later snapshot, \textbf{+29\%} (median +22\% per game);
GTA~V on PS4 went from 12.6 to 19.4 M. The 2015--2016 figures of the main file are
therefore under-estimated, and the back catalogue has real value.

% =====================================================================
\section{The Power BI dashboard}\label{sec:dash}

The dashboard is delivered as the template \nolinkurl{powerbi/Video_Game_Analytics_Dashboard.pbit}.
On opening, Power BI asks for \texttt{DataFolder}, the folder \texttt{data/processed/}, and
loads the four cleaned tables, so every number matches the notebook. The model has five
tables (\texttt{Games}, \texttt{Regional Sales}, \texttt{PS4 XOne}, \texttt{PS4 XOne Regions},
\texttt{Cleaning Log}), one relationship (\texttt{Regional Sales} $\to$ \texttt{Games} on
\texttt{Game ID}) and 22 DAX measures. \emph{Top Genre}, \emph{Top Platform} and \emph{Top
Publisher} are computed with \texttt{TOPN} and follow the filters.

\begin{table}[H]\centering\small
\caption{Pages of the dashboard.}
\begin{tabularx}{\linewidth}{@{}lX@{}}
\toprule
Page & Visuals\\\midrule
1. Overview & KPI cards (Total Games, Total Releases, Total Sales, Top Genre, Top Platform, Top Publisher); sales trend; region donut; sales by genre, top 10 platforms and publishers; scrolling best sellers\\
2. Genres \& Platforms & releases by genre; average and median units per release by genre; units per year by manufacturer; top 12 platforms\\
3. Publishers \& Games & top 10 publishers; releases vs units per release by publisher; top 10 releases; top 10 games, all platforms\\
4. Regions & region donut; region share by year; genre share inside each region; units by region\\
5. Ratings \& Sales & score cards; median units by critic band; units at each critic score; critics vs players; ESRB rating\\
6. PS4 \& Xbox One & cards of the console files; region donut; top 10 games; genres; units by year and console\\
7. Insights & the insights and recommendations of Section~\ref{sec:ins}\\
8. Data Quality & pipeline summary and cleaning log\\
\bottomrule
\end{tabularx}
\end{table}

\textbf{Design.} The report uses a \emph{Neon Arcade} theme made for the gaming audience:
an illustrated night background with a synthwave grid, a header banner with the logo and a
pill-shaped page navigator, KPI cards with icons, translucent rounded panels behind every
visual and one neon colour per chart (cyan, magenta, violet, amber, lime). The palette is
kept consistent: the four regions always have the same colours.

\textbf{Interactivity.} Four slicers synchronised across the analysis pages (year range,
genre, manufacturer, platform); cross-filtering between visuals (a click on a bar filters
the page); Top~N filters that follow the slicers; the page navigator; tooltips on every
visual. The build guide (\nolinkurl{powerbi/Dashboard_Build_Guide.md}) lists the values to
check after opening; every capture below, taken with no filter, shows these values.

\subsection{Overview}
\capture{dashboard/p1_overview.png}{\label{fig:p1}Overview. The six KPI cards answer the brief at a glance:
11{,}563 games, 16{,}715 releases, 8{,}917.9 million units, and the three leaders (Action, PlayStation~2,
Nintendo). The area chart shows the rise to the 2008 peak and the decline; the donut gives the
regional split (North America 49.4\%); the three bar charts rank genres, platforms and publishers;
the ticker at the bottom scrolls the best-selling games (Wii Sports 82.5 M, GTA~V 56.6 M).}
The Overview is the one-page summary required by the brief. Every card is a DAX measure that
follows the slicers: selecting the genre \emph{Role-Playing}, for example, makes Nintendo DS the top
platform, and selecting the manufacturer \emph{Nintendo} makes the Wii the top platform.

\subsection{Genres and platforms (Q1, Q2, Q5)}
\capture{dashboard/p2_genres_platforms.png}{\label{fig:p2}Genres \& Platforms. Left to right: the most
common genres (Action 3.4~K releases), the average and the median units per release by genre (Platform
first in both: 0.93~M and 0.27~M), the stacked units per year by manufacturer (Sony and Nintendo
dominate the 2000s peak) and the twelve best-selling platforms (PS2 1{,}255.6~M).}
Putting frequency next to the average and the median makes the key nuance of Q1 visible:
Action is the most common genre, but Platform and Shooter games sell the most per release,
and the medians are three to four times below the means.

\subsection{Publishers and games (Q3, Q4)}
\capture{dashboard/p3_publishers_games.png}{\label{fig:p3}Publishers \& Games. The top 10 publishers
(Nintendo 1{,}788.8~M, EA 1{,}117.0~M); the scatter of releases against units per release separates the
\emph{hit} model (Nintendo, top left: few releases, 2.5~M each) from the \emph{volume} model (Electronic
Arts, right: 1{,}355 releases); the table of the best-selling releases and the top 10 games summed
over platforms (GTA~V second with 56.6~M).}

\subsection{Regions (Q6)}
\capture{dashboard/p4_regions.png}{\label{fig:p4}Regions. Share of units by region (North America 49.37\%,
Europe 27.2\%, Japan 14.56\%, Other 8.88\%); share of each year's units (Japan's weight collapses after
1995, Europe grows to the first rank in 2016); genre share inside each region (Japan's role-playing
peak and shooter gap) and units by region.}

\subsection{Ratings and sales (Q7)}
\capture{dashboard/p5_ratings_sales.png}{\label{fig:p5}Ratings \& Sales. Cards: 8{,}135 scored releases,
average critic score 69.0, average user score 71.3 (x10), median 0.24~M units. The median per critic band
rises from 0.12~M (below 50) to 1.54~M (90+); the average units at each score bend upwards above 80; critics
and players agree only loosely (wide cloud around the trend line); by ESRB rating, Mature games lead
among the large groups.}
Two readings need care on this page. The \emph{No score} band (0.12~M) gathers the releases
never reviewed by Metacritic, mostly older or smaller games. In the ESRB chart, \textbf{AO
(1.95~M) is a single release} and \textbf{EC} has only eight: the meaningful comparison is between
M (0.31~M), E10+, E and T (about 0.20~M) and \emph{Not rated} (0.11~M).

\subsection{After 2016: PS4 and Xbox One}
\capture{dashboard/p6_ps4_xbox_one.png}{\label{fig:p6}PS4 \& Xbox One (later VGChartz snapshot). Cards:
1{,}135 releases with sales, 864.7~M units, 0.76~M per release, Japan 4.1\% of units. The donut shows
North America 43.3\% and Europe 39.1\%; GTA~V leads with 28.1~M units on the two consoles; shooters lead
the genres (227.2~M); the PS4 outsells the Xbox One every year.}

\subsection{Insights and data quality}
\capture{dashboard/p7_insights.png}{\label{fig:p7}Insights: the nine findings of the analysis and the
recommendations, each tile with its key figures.}
\capture{dashboard/p8_data_quality.png}{\label{fig:p8}Data Quality: the five steps from the raw files to
the model, and the cleaning log loaded from \texttt{cleaning\_log.csv} (16{,}719 rows to 16{,}715 releases,
99.97\% of units kept).}
The last page makes the cleaning auditable inside the dashboard itself: a reader can see
which rows were removed, merged or recoded, and on which base each visual is computed.

% =====================================================================
\section{Insights and recommendations}\label{sec:ins}

\begin{enumerate}
\item \textbf{Choose genres by return per release, not by popularity.} Platform, Shooter,
Sports and Racing have the highest medians; crowded genres with low medians (Adventure,
Strategy, Puzzle) need a clear differentiation or a small budget.
\item \textbf{Invest in quality as judged by the press.} Median sales jump above a critic
score of 80 (0.59 M) and 90 (1.54 M): polish, review-ready launches and press relations
pay; user scores are a weaker signal.
\item \textbf{Release on several platforms.} GTA~V and Call of Duty reach the top by adding
platforms and generations; plan ports early for the leading platforms of the current
generation.
\item \textbf{Adapt the offer to each region.} Shooters and sports for North America and
Europe; role-playing games, handheld and mobile formats for Japan; give Europe the same
launch effort as North America.
\item \textbf{Manage a portfolio of bets.} Half of the market is made by 6.6\% of releases:
plan with medians and hit probabilities, fund a few potential blockbusters, keep small
titles cheap.
\item \textbf{Exploit the back catalogue}: established games still sell 29\% more within
about two years (price cuts, bundles, remasters, ports).
\item \textbf{Complete the data with digital sales} before deciding on the future of a genre
or platform: retail units no longer describe the whole market.
\end{enumerate}

% =====================================================================
\section{Limitations}
\begin{itemize}
\item Retail units only: no prices, revenue, digital, mobile or free-to-play sales.
\item VGChartz figures are estimates and are revised; 2016 is incomplete.
\item Scores exist for half of the releases, almost none before 1996.
\item The correlations are associations, not causes.
\end{itemize}

\section{Reproducibility}

Every number, table and figure of this report is produced by the notebook
\nolinkurl{Video-Game-Analytics.ipynb} from the three files of \nolinkurl{data/raw/}. The
notebook writes the cleaned dataset and the dashboard tables (\nolinkurl{data/processed/}),
the result tables (\nolinkurl{results/}) and the figures (\nolinkurl{figures/}). The template
\nolinkurl{powerbi/Video_Game_Analytics_Dashboard.pbit} loads \nolinkurl{data/processed/}.
The project is published at \url{https://github.com/koffifidelek59-collab/video-game-analytics}.

\section*{References}
{\small\begin{enumerate}[itemsep=0pt]
\item Kirubi, R. (2016). \emph{Video Game Sales with Ratings}, Kaggle (VGChartz sales, Metacritic scores, extraction of 22 December 2016). \url{https://www.kaggle.com/datasets/rush4ratio/video-game-sales-with-ratings}
\item VGChartz. Game sales database (PS4 and Xbox One files). \url{https://www.vgchartz.com}
\item Metacritic. Critic and user scores. \url{https://www.metacritic.com}
\item Entertainment Software Rating Board. Ratings guide. \url{https://www.esrb.org/ratings-guide/}
\item Gebru, T. et al. (2021). Datasheets for Datasets. \emph{Communications of the ACM}, 64(12), 86--92.
\item Microsoft. Power BI documentation. \url{https://learn.microsoft.com/power-bi/}
\end{enumerate}}

\enlargethispage{4\baselineskip}
\vspace{0.1cm}
\noindent\hfill\begin{minipage}{7cm}\begin{flushright}\singlespacing
\textbf{KOUAME Koffi Fid\`ele}\\
Data Analysis Internship\\
\href{mailto:koffifidelek59@gmail.com}{koffifidelek59@gmail.com}
\end{flushright}\end{minipage}

\end{document}
